% Generated draft: requires mathematical and source review.
\documentclass[11pt]{article}
\usepackage[margin=25mm]{geometry}
\usepackage{fontspec}
\setmainfont{Latin Modern Roman}
\usepackage{amsmath,amssymb,mathtools,unicode-math}
\setmathfont{Latin Modern Math}
\usepackage{enumitem,needspace,xurl,hyperref}
\setlength{\parindent}{0pt}
\setlength{\parskip}{5pt}
\setlength{\emergencystretch}{3em}
\newcommand{\N}{\mathbb{N}}
\newcommand{\Z}{\mathbb{Z}}
\newcommand{\Q}{\mathbb{Q}}
\newcommand{\R}{\mathbb{R}}
\newcommand{\C}{\mathbb{C}}
\newcommand{\F}{\mathbb{F}}
\newcommand{\A}{\mathbb{A}}
\newcommand{\PP}{\mathbb{P}}
\newcommand{\E}{\mathbb{E}}
\newcommand{\eps}{\varepsilon}
\newcommand{\dd}{\,\mathrm{d}}
\begin{document}
\section*{UnsolvedMath ID: 30005693}
\textbf{Source problem code:} catalogue entry
\subsection{Definitions and mathematical statement}
Let \(E/F\) be a field extension such that \(\operatorname{char}(F)=0\) and \(\operatorname{trdeg}_F(E)<\infty\). No positive lower bound on the transcendence degree is supplied. Define
\[
G=\operatorname{Aut}_F(E)
 =\{g:E\longrightarrow E\mid g\text{ is a field automorphism and }g|_F=\mathrm{id}_F\}.
\]
The multiplicative abelian group \(E^\times=E\setminus\{0\}\) has the natural \(G\)-action \(g\cdot a=g(a)\).

The supplied problem asks: determine the discrete first cohomology group
\[
 dH^1(G,E^\times).
\]
The supplied record does not explicitly specify the topology on \(G\) or the coefficient category and convention defining \(dH^1\). These qualifications must be resolved before the target is fully specified.

For clarity, one possible interpretation is the following, conditional on the intended conventions. Endow \(G\) with the topology of pointwise convergence on the discrete set \(E\), so that the subgroups fixing finite subsets of \(E\) pointwise form a neighbourhood basis at the identity, and give \(E^\times\) the discrete topology. Its natural action is then continuous. If \(dH^1\) means continuous first group cohomology with discrete abelian coefficients, its multiplicative cocycle description is
\[
 Z^1_{\mathrm{cont}}(G,E^\times)
 =\{c:G\longrightarrow E^\times\mid c\text{ is continuous and }
 c(gh)=c(g)\,g(c(h))\text{ for all }g,h\in G\},
\]
\[
 B^1(G,E^\times)
 =\{c_a\mid a\in E^\times,\quad c_a(g)=a^{-1}g(a)\text{ for all }g\in G\},
 \qquad
 dH^1(G,E^\times)=Z^1_{\mathrm{cont}}(G,E^\times)/B^1(G,E^\times).
\]
This conditional description identifies a conventional target; it is not a computation or confirmation of the source's intended definition. In particular, no algebraic closedness hypotheses on either field, no local compactness hypothesis on \(G\), and no rational vector-space structure on \(E^\times\) are supplied, and none is imposed here.
\subsection{Short English statement}
For a characteristic-zero field extension \(E/F\) of finite transcendence degree, determine \(dH^1(\operatorname{Aut}_F(E),E^\times)\), with the natural automorphism action on nonzero field elements. The topology and precise discrete-cohomology convention require confirmation. The supplied registry labels the problem open; this draft does not independently verify that status.
\subsection{Sources}
\begin{itemize}
\item Supplied local UnsolvedMath registry record, permanent ID \(30005693\), Oberwolfach Reports Open Problems collection, exact source problem code \texttt{catalogue entry}. The record supplies both the field-extension formulation and an original excerpt asking for \(dH^1(G,E^\times)\). Its labels \texttt{corrected\_verified} and \texttt{open}, and its claim of checking the official report, are supplied provenance assertions, not independent verification by this draft.
\item Supplied report citation: \emph{Mini-Workshop: Homological aspects of TDLC-groups (2024)}, Oberwolfach Reports, DOI 10.4171/owr/2023/55, pp.\ 3103--3138. \url{https://doi.org/10.4171/owr/2023/55}. The supplied report link is also \url{https://ems.press/content/serial-article-files/48166}.
\item Attribution in the original excerpt: Thomas Weigel, \emph{Discrete cohomology, the Hattori-Stallings rank, the Euler characteristic and formal Dirichlet series for t.d.l.c. groups}, joint work with Ilaria Castellano, Gianmarco Chinello, Bianca Marchionna and George Willis. The registry does not separately identify a proposer.
\item Additional supplied literature reference, labelled ``Castellano'': \url{https://doi.org/10.1017/S0305004118000762}. No title or fuller bibliographic information is supplied for this reference.
\item Literature assessment dated exactly \texttt{2026-08-22}: the supplied curation labels the problem open and reports that the source leaves this specialized calculation unanswered, while later general theory of rational discrete first cohomology does not compute this coefficient case. This is supplied dated curation, not independently verified literature assessment; its status and cited claims require independent checking.
\end{itemize}
\end{document}
